% Minimal EMNLP template smoke test.
\pdfoutput=1
\documentclass[11pt]{article}

% Remove [review] for camera-ready; for smoke test either is fine.
\usepackage[review]{EMNLP2023}

\usepackage{times}
\usepackage{latexsym}
\usepackage[T1]{fontenc}
\usepackage[utf8]{inputenc}
\usepackage{microtype}
\usepackage{inconsolata}

\usepackage{graphicx}
\usepackage{booktabs}
\usepackage{hyperref}

\title{GALILEO: Measuring Multi-turn Belief-Consistency Failures and Recovery Under Social Pressure}
\author{Anonymous Authors}
\date{}

\begin{document}
\maketitle

\begin{abstract}
Large language models (LLMs) can change previously correct answers when confronted with social pressure such as repeated denial, authority claims, or persuasive reframing.
Prior work has documented sycophancy and persuasion effects, but there is no widely adopted, end-to-end protocol for ground-truth tasks that (i) measures \emph{when} the first failure occurs (turn-of-failure), (ii) summarizes degradation over multiple rounds (survival curves), and (iii) evaluates whether models can return to the original truth after a flip (recovery).

We introduce \textbf{GALILEO}, a reproducible evaluation pipeline for ground-truth tasks (math, extractive QA, multiple-choice QA, and open-domain QA).
Given an initially correct answer, GALILEO applies five adversarial personas (soft pressure, simple denial, strong pressure, authority claim, and logical trap) for up to five rounds, then prompts for recovery on flipped samples.
Across multiple model families and seeds, we find that persona pressure consistently worsens survival, increases fail@1, and shifts turn-of-failure earlier compared to a non-adversarial drift baseline (neutral re-asking control).
We further observe that \emph{robustness} (staying correct) and \emph{recoverability} (returning to truth after a flip) behave as distinct axes.
\end{abstract}

\section{Introduction (placeholder)}
\label{sec:intro}

\noindent\textbf{Contributions.} GALILEO provides:
(1) a ground-truth, multi-turn protocol with survival and turn-of-failure metrics;
(2) a drift baseline (neutral re-asking control) that separates generic answer instability from persona-induced effects; and
(3) a recovery stage that quantifies return-to-truth after a flip.

\section{Related Work}
\label{sec:related}

Prior work has documented sycophancy and social-pressure failure modes in instruction-tuned LLMs, including analyses of when models defer to user assertions even when those assertions conflict with task truth \citep{sharma2023understandingsycophancylanguagemodels}. Subsequent benchmarks and diagnostics extend this agenda to broader settings and measurement choices, e.g., sycophancy-focused evaluation suites \citep{fanous2025sycevalevaluatingllmsycophancy,cheng2025elephantmeasuringunderstandingsocial} and multi-agent/peer-pressure scenarios \citep{song2025kairospeerpressure}.

A complementary thread studies stability and consistency under sequential interaction, highlighting that models may change answers across turns even without explicit adversarial pressure \citep{li2025firmorfickle,tosato2025persistentinstabilityllmspersonality}. GALILEO builds on these observations by (i) grounding evaluation in tasks with externally checkable correctness, (ii) explicitly separating generic multi-turn drift (neutral re-asking control) from persona-induced pressure, and (iii) adding a recovery stage to quantify return-to-truth after a flip. Related work also examines belief resistance under strategic persuasion in longer conversations \citep{huang2026vulnerabilityllmsbeliefsystems}; our protocol targets a more standardized, end-to-end measurement stack with survival/turn-of-failure metrics and reproducible exports.

More broadly, our emphasis on verifiable truth connects to work on truthfulness benchmarks such as TruthfulQA \citep{lin2022truthfulqameasuringmodelsmimic}, while focusing on a different axis: multi-turn belief-consistency failures and recovery under conversational pressure.

\section{Figure include smoke test}

\begin{figure}[t]
  \centering
  \includegraphics[width=\linewidth]{figures/protocol_overview}
  \caption{Protocol overview (PDF include smoke test).}
\end{figure}

\begin{figure}[t]
  \centering
  \includegraphics[width=\linewidth]{figures/cross_family_survival_r5_control_vs_logicaltrap_seed1-2_20260217}
  \caption{Cross-family survival (PDF include smoke test).}
\end{figure}

\begin{figure}[t]
  \centering
  \includegraphics[width=\linewidth]{figures/decoding_sweep_qwen_delta_seed1-2_20260211}
  \caption{Appendix robustness: decoding sensitivity (PDF include smoke test).}
  \label{fig:decoding-sweep}
\end{figure}

\clearpage
\section{Claim\texorpdfstring{$\rightarrow$}{->}evidence table smoke test}

\begin{table}[t]
\centering
\small
\begin{tabular}{p{0.14\linewidth} p{0.30\linewidth} p{0.30\linewidth} p{0.20\linewidth}}
\toprule
Claim & Evidence & Tracked artifacts & Reproducer \\
\midrule
C1 & Survival curves + TOF & survival\_curve\_personawise...; tof\_personawise... & validate\_paper\_exports.py (results\_paper) \\
C2 & Table W deltas (control vs persona) & table\_w\_effect\_delta... & make\_table\_w... + validate \\
C6 & Fig.~\ref{fig:decoding-sweep} & decoding\_sweep\_qwen\_temp\_summary... & results\_paper/qwen\_temp* \\
\bottomrule
\end{tabular}
\caption{Minimal claim-to-evidence table (format smoke test).}
\end{table}

\clearpage
\bibliographystyle{acl_natbib}
% Use the repo-wide BibTeX database at the repo root.
% (This keeps citations in sync with the writing scaffold in docs/paper/.)
\bibliography{../../../references}

\end{document}
